\documentclass[11pt]{article}
\usepackage{amsmath,amssymb,amsfonts}
\usepackage[margin=1in]{geometry}
\usepackage{bm}
\usepackage{hyperref}

\newcommand{\uu}{\bm{u}}
\newcommand{\vv}{\bm{v}}
\newcommand{\ww}{\bm{w}}
\newcommand{\ff}{\bm{f}}
\newcommand{\nn}{\bm{n}}
\newcommand{\xx}{\bm{x}}
\newcommand{\grad}{\nabla}
\newcommand{\dive}{\operatorname{div}}
\newcommand{\Reyn}{\mathrm{Re}}
\newcommand{\Sten}{\bm{S}}
\newcommand{\dx}{\,d\xx}

\title{The Cocquet et al.\ Darcy--Brinkman--Forchheimer formulation\\[2pt]
\large (reference note for the \texttt{CocquetTubeTest/data/structured} validation)}
\author{}
\date{}

\begin{document}
\maketitle

\noindent
This note records, for future reference, the \emph{exact} continuous and discrete formulation of
\begin{quote}
P.-H.\ Cocquet, M.\ Rakotobe, D.\ Ramalingom, A.\ Bastide,
\emph{Error analysis for the finite element approximation of the Brinkman--Darcy--Forchheimer model
for porous media with mixed boundary conditions}, J.\ Comput.\ Appl.\ Math.\ (2020),
\href{https://hal.univ-reunion.fr/hal-02561058}{hal-02561058}.
\end{quote}
The PDF is in \texttt{theory/Cocquet et al.\ - 2021 - Error analysis\ldots pdf}; equation numbers below
are the paper's. The last section maps every term onto our code so the
\texttt{CocquetTubeTest/data/structured} ``Galerkin'' comparison run reproduces the paper's method exactly.

\section{Continuous problem}

On a bounded Lipschitz domain $\Omega\subset\mathbb{R}^d$ with spatially varying porosity
$\varepsilon:\Omega\to(0,1]$, the dimensionless steady Darcy--Brinkman--Forchheimer (DBF) model
reads (Eq.~1):
\begin{equation}
\left\{
\begin{aligned}
-\dive\!\big(2\Reyn^{-1}\varepsilon\,\Sten(\uu) - \varepsilon\,\uu\otimes\uu\big)
   + \varepsilon\grad p + \alpha(\varepsilon)\uu + \beta(\varepsilon)|\uu|\,\uu
   &= \varepsilon\ff, && \text{in }\Omega,\\
\dive(\varepsilon\uu) &= 0, && \text{in }\Omega,
\end{aligned}
\right.
\end{equation}
where $\Sten(\uu)=\tfrac12(\grad\uu+\grad\uu^{\top})$ is the symmetric velocity gradient, $\Reyn$ the
Reynolds number, and $\ff$ an external force. One always has $\alpha(1)=\beta(1)=0$, so the standard
Navier--Stokes equations are recovered when $\varepsilon\equiv1$. Using
$\dive(\varepsilon\uu\otimes\uu)=\dive(\varepsilon\uu)\uu+\varepsilon(\uu\cdot\grad)\uu$ and
$\dive(\varepsilon\uu)=0$, the convective term is $\varepsilon(\uu\cdot\grad)\uu$.

\paragraph{Mixed boundary conditions (Eq.~2).} With $\partial\Omega=\Gamma_w\cup\Gamma_{\mathrm{in}}\cup\Gamma_{\mathrm{out}}$,
\begin{equation}
\uu = \bm{0}\ \text{on }\Gamma_w,\qquad
\uu = \uu_{\mathrm{in}}\ \text{on }\Gamma_{\mathrm{in}},\qquad
\varepsilon\big(2\Reyn^{-1}\Sten(\uu)-p\big)\nn = \bm{0}\ \text{on }\Gamma_{\mathrm{out}},
\end{equation}
i.e.\ no-slip walls, an inflow Dirichlet datum, and a traction-free (pseudo-traction) Neumann outlet.
The wall/outlet corners are Dirichlet--Neumann junctions; their corner singularity is the subject of
the convergence study (see \texttt{docs/cocquet\_convergence\_analysis.md}).

\paragraph{Darcy and Forchheimer closure (Eq.~49).} For packed beds,
\begin{equation}
\alpha(\varepsilon)=\frac{150}{\Reyn}\Big(\frac{1-\varepsilon}{\varepsilon}\Big)^{2},
\qquad
\beta(\varepsilon)=1.75\,\frac{1-\varepsilon}{\varepsilon}.
\end{equation}

\section{Weak (Galerkin) formulation}

With $\bm{X}=\{\vv\in H^1(\Omega)^d:\vv|_{\Gamma_w\cup\Gamma_{\mathrm{in}}}=\bm 0\}$ and
$\bm{F}=\varepsilon\ff$, the variational problem is (Eq.~4): find
$(\uu,p)\in \bm{X}_1\times L^2(\Omega)$ with $\uu|_{\Gamma_{\mathrm{in}}}=\uu_{\mathrm{in}}$ such that
\begin{equation}
\left\{
\begin{aligned}
a(\varepsilon;\uu,\vv) + c(\varepsilon;\uu,\uu,\vv) + b(\varepsilon;\vv,p) &= \langle\bm F,\vv\rangle,
  && \forall\,\vv\in\bm X,\\
b(\varepsilon;\uu,q) &= 0, && \forall\,q\in L^2(\Omega),
\end{aligned}
\right.
\end{equation}
\begin{align}
a(\varepsilon;\uu,\vv) &= 2\Reyn^{-1}\!\int_\Omega \varepsilon\,\Sten(\uu){:}\Sten(\vv)\dx
   + \int_\Omega \alpha(\varepsilon)\,\uu\cdot\vv\dx,\\
c(\varepsilon;\uu,\vv,\ww) &= \int_\Omega \varepsilon\,(\uu\cdot\grad)\vv\cdot\ww\dx
   + \int_\Omega \beta(\varepsilon)\,|\uu|\,\vv\cdot\ww\dx,\\
b(\varepsilon;\vv,q) &= -\int_\Omega q\,\dive(\varepsilon\vv)\dx .
\end{align}
The porosity weights \emph{every} operator (viscous, convective, pressure, divergence, source); the
reaction coefficients $\alpha,\beta$ depend on $\varepsilon$ pointwise.

\section{Discretization (the method we reproduce)}

\begin{itemize}
\item \textbf{Taylor--Hood $P_2/P_1$}, inf-sup stable, so \textbf{no stabilization} is used --- the
discrete problem is the pure Galerkin restriction of Eq.~4 to $\bm X_h\times M_h$ (Eq.~29/40).
\item A small pressure penalty $\eta\,p$ ($\eta=10^{-7}$) is added to the incompressibility equation so
every sub-matrix is invertible for the direct solver.
\item The porosity is a \textbf{$P_1$ finite-element interpolant} $\varepsilon_h$ of $\varepsilon$,
regardless of the velocity order.
\item The nonlinear discrete system is solved by a \textbf{Picard fixed-point iteration} (Eq.~41):
freeze the convecting velocity and the Forchheimer factor $|\uu_{h,n-1}|$, solve the resulting linear
saddle-point problem for $(\uu_{h,n},p_{h,n})$, repeat. (The discrete solution is independent of the
nonlinear solver; our driver uses an exact-Newton/Picard cascade and reaches the same root.)
\end{itemize}

\paragraph{Numerical test case (\S4.2).} Smooth porosity
$\varepsilon(x,y)=0.45\big(1+\tfrac{1-0.45}{0.45}\exp(-(1-y))\big)=0.45+0.55\,e^{\,y-1}$,
domain $\Omega=(0,2)\times(0,1)$, $\Gamma_{\mathrm{in}}=\{0\}\times[0,1]$,
$\Gamma_{\mathrm{out}}=\{2\}\times[0,1]$, $\Gamma_w=[0,2]\times(\{0\}\cup\{1\})$, parabolic inflow
$\uu_{\mathrm{in}}(y)=c_{\mathrm{in}}\,y(1-y)\,\bm e_x$. The convergence figures (Figs.~2--3) use a
reference solution at $N=200$ and coarse meshes $N\le100$, with $h=\sqrt2/N$ (triangle diameter).
The paper observes optimal $\mathrm{Err}_{tot}=O(h^2)$ in the energy norm; the velocity $L^2$ error does
\emph{not} gain the extra order (it stays $O(h^2)$), which the paper attributes to the discrete
porosity and we attribute (for our equal-order method) to the outlet corner singularity.

\section{Mapping to the code}

The production solver already implements Eq.~4 term-for-term; the only thing the Cocquet comparison
needs beyond our VMS runs is to switch the stabilization off. No \texttt{src/} change is required.

\begin{center}
\begin{tabular}{@{}ll@{}}
\hline
Cocquet term & Code (\texttt{src/}) \\
\hline
$2\Reyn^{-1}\varepsilon\,\Sten(\uu){:}\Sten(\vv)$ & \texttt{SymmetricGradientViscosity}, $\nu=\Reyn^{-1}$ \\
$\alpha(\varepsilon)\uu+\beta(\varepsilon)|\uu|\uu$, Eq.~49 & \texttt{ForchheimerErgunLaw(sigma\_linear,\,sigma\_nonlinear)}; \\
 & \texttt{sigma\_linear}$=150/\Reyn$, \texttt{sigma\_nonlinear}$=1.75$ \\
$\varepsilon(\uu\cdot\grad)\uu$, $\varepsilon\grad p$, $\dive(\varepsilon\uu)$ & porosity-weighted base terms in \\
 & \texttt{build\_stabilized\_weak\_form\_residual} (\texttt{continuous\_problem.jl}) \\
$\eta\,p$ penalty & \texttt{eps\_val} ($=\eta$) in the mass equation \\
$P_1$ porosity $\varepsilon_h$ & \texttt{porosity\_order=1} in \texttt{run\_convergence.jl::build\_solver} \\
Taylor--Hood $P_2/P_1$ & element pair \texttt{(k\_velocity,k\_pressure)=(2,1)} \\
\textbf{no stabilization} (pure Galerkin) & weak-form builders called with
   \texttt{mult\_mom=mult\_mass=0} \\
\hline
\end{tabular}
\end{center}

\paragraph{Where the Galerkin run lives.} The unstabilized assembly is driven by
\texttt{test/extended/CocquetTubeTest/galerkin\_driver.jl} (\emph{not} \texttt{src/}: it is validation
tooling). It reuses the production builders
\texttt{build\_stabilized\_weak\_form\_residual} / \texttt{build\_picard\_jacobian} /
\texttt{build\_stabilized\_weak\_form\_jacobian} with \texttt{mult\_mom=mult\_mass=0}, which zero the VMS
stabilization terms ($\bm L^\ast\!\cdot\tau\,\mathcal R$) and leave exactly the bare Galerkin DBF system
of Eq.~4 (with the $\eta p$ penalty). The \texttt{comparison\_runs} list in
\texttt{test/extended/CocquetTubeTest/data/structured/paper\_comparison.json} runs our VMS $P_1/P_1$ and $P_2/P_2$ (method \texttt{ASGS}) alongside
this Cocquet $P_2/P_1$ Galerkin run (method \texttt{Galerkin}, $P_1$ porosity), and
\texttt{plot\_convergence.py} draws all three in one figure.

\section{Mesh generation and the paper's $h=\sqrt2/N$ inconsistency}
\label{sec:mesh-generation}

Reproducing Cocquet's Figure~2 quantitatively requires reproducing their meshes. The paper specifies
the meshes verbatim in \S4.2 (p.~30):
\begin{quote}\itshape
``the mesh is obtained thanks to the Freefem command \texttt{buildmesh(a(N)+b(N)+c(N)+d(N))} with $N$
being the number of vertices on each part of the boundary denoted by $a,b,c,d$. As a result, the mesh-size
$h$ is $h=\sqrt2/N$, and we can consider only $N$.''
\end{quote}

\paragraph{What FreeFem actually does with that command.} In FreeFem++, the integer argument of a border
specifies the number of \emph{segments} along that border. Concretely, with the four borders
\begin{align*}
a(t\in[0,2])\!: &\quad x=t,\ y=0\quad(\text{bottom wall, length }2)\\
b(t\in[0,1])\!: &\quad x=2,\ y=t\quad(\text{outlet, length }1)\\
c(t\in[2,0])\!: &\quad x=t,\ y=1\quad(\text{top wall, length }2)\\
d(t\in[1,0])\!: &\quad x=0,\ y=t\quad(\text{inlet, length }1),
\end{align*}
the call \texttt{buildmesh(a(N)+b(N)+c(N)+d(N))} discretises each border into exactly $N$ equal segments
\emph{regardless of its physical length}. Wall edges (length $2$) therefore have length $2/N$, and
inlet/outlet edges have length $1/N$ --- the boundary, and hence the interior, is \textbf{anisotropic}.

\paragraph{Why the paper's $h=\sqrt2/N$ is inconsistent with the literal recipe.} The value $h=\sqrt2/N$
is the diameter (hypotenuse) of an \emph{isotropic} right triangle whose two legs are both $1/N$. With
the literal $a(N)+b(N)+c(N)+d(N)$ recipe the triangles have legs $\sim1/N$ and $\sim2/N$, so the diameter
is $\sqrt{(1/N)^2+(2/N)^2}=\sqrt5/N\approx2.24/N$, \emph{not} $\sqrt2/N$. The two statements
``$N$ vertices on each border'' and ``$h=\sqrt2/N$'' cannot both be true. Empirically the literal recipe
produces $\sim N\!\times\!N\!\times2$ cells (e.g.\ $170$ at $N=10$, $63\,872$ at $N=200$); a mesh truly
with $h=\sqrt2/N$ would have $\sim2N\!\times\!N\!\times2\approx 4N^2$ cells (e.g.\ $400$ at $N=10$,
$\sim160\,000$ at $N=200$) --- a factor of $\sim 2.5\!\times$ in cell count and a corresponding factor
of $\sim 2$--$3\times$ in absolute discretisation error.

\paragraph{Resolution: uniform-edge correction.} We resolve the inconsistency in favour of the
\textbf{stated mesh size} $h=\sqrt2/N$, on the grounds that (i)~$h$ is the quantity the paper plots
against and reports convergence rates in, and (ii)~uniform edges $1/N$ everywhere are the natural way to
``consider only $N$'' (the paper's phrase) without separately tracking the two different boundary
spacings. The corresponding FreeFem command is
\[
  \boxed{\ \texttt{buildmesh(a(2*N) + b(N) + c(2*N) + d(N))}\ }
\]
i.e.\ \textbf{$2N$ segments on the length-$2$ walls} and $N$ on the length-$1$ inlet/outlet, so every
boundary edge has length $1/N$ and the interior --- still meshed by FreeFem's Delaunay/BAMG generator
without any transfinite surface constraint --- inherits roughly the same density. The triangle diameter
is then genuinely $\sqrt2/N$, matching the paper's $h$.

\paragraph{Cell counts and code wiring.} The corrected meshes are produced by
\texttt{test/extended/CocquetTubeTest/data/freefem\_meshes/meshes/create\_paper\_meshes.edp}
(FreeFem run separately --- not built into
our Julia toolchain because no macOS-arm64 FreeFem binary exists) and saved as the FreeFem-native
\texttt{.msh} format. \texttt{src/run\_simulation.jl::load\_freefem\_mesh} translates a FreeFem mesh into
a Gridap \texttt{DiscreteModel} via an in-line conversion to Gmsh 2.2 format, mapping the FreeFem
boundary labels $1,2,3,4$ onto named physical groups \texttt{"inlet"}, \texttt{"outlet"}, \texttt{"walls"}
(both walls share one named group) and a 2D \texttt{"domain"} group; these names then resolve directly in
\texttt{build\_fe\_spaces(\ldots,\ ["inlet","walls"],\ \ldots)}. The
\texttt{CocquetTubeTest/data/freefem\_meshes} driver opts into the FreeFem path via the
\texttt{freefem\_mesh\_dir} key in its config; the file naming convention is
\texttt{mesh\_fig2\_N<N>.msh} for coarse meshes and \texttt{mesh\_fig2\_N<N>\_reference.msh} for the
$N=200$ reference. Indicative cell counts:
\begin{center}
\small
\begin{tabular}{@{}lrrr@{}}
\hline
$N$ & \emph{literal} $a(N){+}b(N){+}c(N){+}d(N)$ & \textbf{uniform} $a(2N){+}b(N){+}c(2N){+}d(N)$ & structured Cartesian \\
\hline
$10$  & $170$    & $448$    & $400$   \\
$40$  & $2\,852$ & $7\,236$ & $6\,400$ \\
$100$ & $\sim17\,800$ & $45\,964$ & $40\,000$ \\
$200$ & $63\,872$ & $178\,460$ & $160\,000$ \\
\hline
\end{tabular}
\end{center}

The uniform recipe brings the FreeFem mesh's cell count into line with both the paper's stated
$h=\sqrt2/N$ resolution and our structured Cartesian benchmark, so the comparison of Galerkin
$P_2/P_1$, VMS $P_1/P_1$ and VMS $P_2/P_2$ on these meshes is now an apples-to-apples reproduction of
Figure~2.

\end{document}
